\chapter{Objetivos}

\section{Objetivo general}

El objetivo general de este trabajo es \textbf{auditar la validez del rendimiento}
de un clasificador binario tumor/no-tumor basado en una red convolucional 3D y
entrenado con resonancia magnética cerebral pública multi-fuente, determinando si
dicho rendimiento refleja la detección de lesión o la explotación de un
\emph{confound} dominio$\leftrightarrow$clase, y \textbf{proponer un protocolo de
auditoría reproducible} que permita diagnosticar este modo de fallo en sistemas
análogos.

El trabajo no persigue, por tanto, entregar una herramienta de triaje
desplegable, sino caracterizar de forma rigurosa el aprendizaje espurio que puede
surgir al combinar datasets públicos en los que la clase coincide con la
procedencia. El clasificador es el objeto de estudio de la auditoría, no su
producto final.

\section{Objetivos específicos}

Del objetivo general se derivan los siguientes objetivos específicos, alineados
con el desarrollo de los capítulos posteriores:

\begin{enumerate}
  \item \textbf{Implementar un pipeline 3D reproducible} para la clasificación de
  RM cerebral multi-fuente, que comprenda un preprocesado homogéneo de todas las
  cohortes, la arquitectura de la CNN 3D y una partición a nivel de sujeto con
  semillas fijas (capítulos~4 y~5).

  \item \textbf{Medir, de forma neutra, el rendimiento aparente} del clasificador
  en el régimen de validación interna sobre el pool multi-fuente, tomándolo como
  punto de partida que ha de auditarse y no como un logro a alcanzar
  (capítulos~6 y~7).

  \item \textbf{Auditar la validez de ese rendimiento} mediante una cadena de
  pruebas anti-confounding: un baseline trivial de intensidad, la validación
  cruzada por dominio (\emph{leave-one-dataset-out}), la validación intra-dominio
  sobre una cohorte de un solo escáner y el análisis del espacio latente
  (capítulos~6 y~7).

  \item \textbf{Descartar explicaciones alternativas} del rendimiento ---la fuga de
  información por partición y la presencia de duplicados--- para aislar el confound
  de dominio como causa (capítulos~6 y~7).

  \item \textbf{Caracterizar el confound dominio$\leftrightarrow$clase} y su
  mecanismo de aprendizaje espurio, explicando por qué se trata de un sesgo
  estructural ligado a la composición de los datos (capítulos~7 y~8).

  \item \textbf{Derivar recomendaciones metodológicas} para el diseño y la
  evaluación de estudios con datos públicos multi-fuente, transferibles a
  problemas análogos (capítulo~8).
\end{enumerate}
